Large language model (LLM) agents are increasingly deployed in mixed-motive settings, but we still lack systematic evidence on how capability gaps shape strategic outcomes and fairness. We study this problem in three controlled negotiation games: item allocation, diplomatic accord negotiation, and participatory budgeting. Across these games, we ask whether higher model capability predicts better outcomes, how competition and cooperation parameters mediate that relationship, whether stronger agents extract above-fair shares, and how these patterns change beyond bilateral play. We benchmark capability with LMArena Elo and evaluate outcomes with both utility and a cross-game exploitation index relative to principled fairness baselines (Nash Bargaining Solution in Games 1--2; Lindahl cost sharing in Game 3). Across all three environments, higher capability generally improves payoff while increasing competition lowers aggregate utility, but the mechanism differs by game. Exploitation signals are strongest in item allocation and diplomatic accords, while participatory budgeting is more coordination-limited under scarcity. Preliminary $N > 2$ experiments suggest capability advantages can persist, and sometimes strengthen, through focal-point and coalition dynamics.
